%% kapa: exam=1 round=20 year=2009 period=2 form=A subject=appraisal_law subject_ko=부동산관계법규 seq=16 number=56
%% source: Q-Net 감정평가사 기출문제 (cst00309) artlSeq=5024178 file=extracted/1008/5024178/2009년 제20회 감정평가사 1차시험 기출문제/2교시.hwp
%% conversion: kordoc 4.24.51 -> tools/gen1.py (auto); review=proofread-sonnet; flags=none
\begin{question}{56}
국토의 계획 및 이용에 관한 법령상 도시계획의 지위에 관한 설명으로 옳은 것은?
\begin{choices}
  \item 다른 법률에 따른 토지의 이용ㆍ개발 및 보전에 관한 계획은 국토의 계획 및 이용에 관한 법률에 의한 도시계획의 기본이 된다.
  \item 특별시장ㆍ광역시장ㆍ시장 또는 군수가 관할 구역에 대하여 다른 법률에 따라 환경ㆍ교통 등에 관하여 수립한 부문별 계획은 도시기본계획에 우선한다.
  \item 중앙행정기관의 장이나 지방자치단체의 장이 다른 법률에 따라 토지 이용에 관한 지역ㆍ지구를 지정하는 경우에는 국토의 계획 및 이용에 관한 법률에 따른 용도지역ㆍ용도지구의 지정목적과 다르게 지정할 수 있다.
  \item 광역도시계획이 수립되어 있는 지역에 대하여 수립하는 도시기본계획의 내용이 광역도시계획과 다를 때에는 도시기본계획의 내용이 우선한다.
  \item 도시관리계획은 광역도시계획과 도시기본계획에 부합되어야 한다.
\end{choices}
\end{question}
\answer{5}
